\chapter*{Conclusion}
\addcontentsline{toc}{chapter}{Conclusion}

This thesis investigated several algorithmic choices for the three-dimensional container packing problem, 
with the goal of understanding how different evolutionary and local-search components influence solution quality.

The experimental results show that selecting placement heuristics evolutionarily instead of relying
on a single fixed heuristic leads to a substantial improvement in performance. 
In particular, allowing the algorithm to choose among several placement heuristics 
significantly outperforms the Max Distance baseline used in the original formulation. 
Similarly, the evolutionary ordering of boxes can often be replaced by simpler heuristic 
orderings based on box volume or box weight without degrading the results. 
In some cases, especially for larger instances, these heuristic orderings even improve convergence by reducing the effective search space.

Another important finding is that a small amount of mutation can be beneficial, 
but mainly for the smallest instances. For larger instances, mutation becomes increasingly 
disruptive and tends to degrade performance. Memetic hill climbing, on the other hand, 
did not bring any practical benefit. Applying a few hill-climbing iterations to offspring 
increased runtime without improving the quality of the solutions. 
When hill climbing was applied to elite individuals, 
the results were even worse, because the search converged faster to local optima.

The influence of box weights was also examined. Adding weights creates a strong additional 
feasibility constraint, which changes the behaviour of the search but does not fundamentally alter the conclusions. 
The elitist genetic algorithm remains clearly superior to hill climbing in all cases. 
In this setting, hill climbing serves as a useful baseline, 
but not as a competitive alternative to the elitist genetic algorithm.

Probabilistic hill climbing was also tested as a baseline inspired 
by simulated annealing. The results indicate that it does not improve 
on standard hill climbing for this problem. In fact, the best configurations 
consistently used low values of the initial acceptance probability 
and the decay coefficient, which means that the method performs best 
when it behaves very similarly to ordinary hill climbing. 
This suggests that accepting worse solutions is not particularly helpful in this search space.

From a software engineering perspective, the thesis also produced a 
fairly extensive and modular codebase implemented in C\#. The program was 
designed so that extending it with additional algorithms or heuristic components should be straightforward. 
At the same time, the implementation grew to roughly 130 source files, which makes the project relatively 
large and somewhat difficult to maintain. In hindsight, a more compact language such as Julia 
might have reduced the amount of boilerplate and simplified experimentation.

Several directions for future work remain open. One interesting extension would be to assign 
different selection probabilities to the individual components of the solution representation. 
At the moment, rotation, box ordering, and placement heuristics are encoded together in a single vector, 
which gives them the same chance of being modified. It would be worthwhile to investigate whether 
better results could be obtained by changing these components at different rates, 
for example by mutating the rotation part more frequently than the ordering or placement part.

Another promising direction is to extend the evolutionary search to the container itself. 
If multiple container dimensions were available, the algorithm could be allowed to choose 
not only box rotations, ordering, and placement heuristics, but also the container variant. 
Such a choice would likely deserve a lower mutation probability than the other components, 
but it could open the way to a more general framework for adaptive packing.

Overall, the thesis demonstrates that the design of the search 
process has a strong impact on solution quality. The best results were obtained by 
combining evolutionary search with flexible heuristic selection, while local-search 
enhancement and probabilistic acceptance did not provide additional benefits for the tested instances.